\begin{thebibliography}{99}

\bibitem{rin63} G. Ringel. \emph{Problem 25}. In: Theory of Graphs and its Applications (Smolenice, 1963), 162.

\bibitem{ros67} A. Rosa. \emph{On certain valuations of the vertices of a graph}. Theory of Graphs (Internat. Sympos., Rome, 1966), Gordon and Breach, 1967, 349--355.

\bibitem{luc1883} E. Lucas. \emph{R\'ecr\'eations math\'ematiques}, tome II. Gauthier-Villars, Paris, 1883.

\bibitem{gal25} J. A. Gallian. \emph{A dynamic survey of graph labeling}. Electron. J. Combin. DS6, 28th ed. (2025).

\bibitem{mps18} R. Montgomery, A. Pokrovskiy, B. Sudakov. \emph{Embedding rainbow trees with applications to graph labelling and decomposition}. J. Eur. Math. Soc. 22 (2020), 3101--3132. \href{https://arxiv.org/abs/1803.03316}{arXiv:1803.03316}

\bibitem{mps20} R. Montgomery, A. Pokrovskiy, B. Sudakov. \emph{A proof of Ringel's conjecture}. Geom. Funct. Anal. 31 (2021), 663--720. \href{https://arxiv.org/abs/2001.02665}{arXiv:2001.02665}

\bibitem{ks20} P. Keevash, K. Staden. \emph{Ringel's tree packing conjecture in quasirandom graphs}. \href{https://arxiv.org/abs/2004.09947}{arXiv:2004.09947}

\bibitem{aagh16} A. Adamaszek, P. Allen, C. Grosu, J. Hladk\'y. \emph{Almost all trees are almost graceful}. Random Structures Algorithms 56 (2020), 948--987. \href{https://arxiv.org/abs/1608.01577}{arXiv:1608.01577}

\bibitem{lpw25} S. Letzter, A. Pokrovskiy, E. Williams. \emph{On the gracesize of trees}. November 2025; revised July 2026. \href{https://arxiv.org/abs/2511.11331v2}{arXiv:2511.11331v2}

\bibitem{dn26} B. Dumitru, M. Nacu. \emph{Alternating Extremes in Graceful Labelings of Full Binary Trees and Spider Trees}. July 2026. \href{https://arxiv.org/abs/2607.12597v2}{arXiv:2607.12597v2}

\bibitem{sz26} S. Shan, Y. Zhong. \emph{Graceful Labeling of Two Families of Spiders}. May 2026. \href{https://arxiv.org/abs/2605.14295v2}{arXiv:2605.14295v2}

\bibitem{gs80} R. L. Graham, N. J. A. Sloane. \emph{On additive bases and harmonious graphs}. SIAM J. Algebraic Discrete Methods 1 (1980), 382--404.

\bibitem{mp25b} A. M\"uyesser, A. Pokrovskiy. \emph{On the Graham--Sloane harmonious labelling conjecture}. \href{https://arxiv.org/abs/2509.05280}{arXiv:2509.05280}

\bibitem{bghj91} R. Beals, J. A. Gallian, P. Headley, D. Jungreis. \emph{Harmonious groups}. J. Combin. Theory Ser. A 56 (1991), 223--238.

\bibitem{mp22} A. M\"uyesser, A. Pokrovskiy. \emph{A random Hall--Paige conjecture}. \href{https://arxiv.org/abs/2204.09666}{arXiv:2204.09666}

\bibitem{jav25} M. Javaheri. \emph{R-harmonious groups}. \href{https://arxiv.org/abs/2512.08830}{arXiv:2512.08830}

\bibitem{jav26} M. Javaheri. \emph{Symmetric harmoniousness of odd-order groups}. Australas. J. Combin. 96(1) (2026), 134--141. \href{https://ajc.maths.uq.edu.au/download/96/ajc_v96_p134.pdf}{Journal PDF (theorem numbering used here)}. Preprint: \href{https://arxiv.org/abs/2511.13430v1}{arXiv:2511.13430v1}

\bibitem{mm84} M. Maamoun, H. Meyniel. \emph{On a problem of G. Hahn about coloured Hamiltonian paths in $K_{2^t}$}. Discrete Math. 51 (1984), 213--214.

\bibitem{hr09} P. Horak, A. Rosa. \emph{On a problem of Marco Buratti}. Electron. J. Combin. 16 (2009), R20.

\bibitem{pp14} A. Pasotti, M. A. Pellegrini. \emph{On the Buratti--Horak--Rosa conjecture about Hamiltonian paths in complete graphs}. \href{https://arxiv.org/abs/1311.2785}{arXiv:1311.2785}

\bibitem{opps20} M. A. Ollis, A. Pasotti, M. A. Pellegrini, J. R. Schmitt. \emph{New methods to attack the Buratti--Horak--Rosa conjecture}. \href{https://arxiv.org/abs/1912.07377}{arXiv:1912.07377}

\bibitem{ao24} O. A\u{g}\i rseven, M. A. Ollis. \emph{Grid-Based Graphs, Linear Realizations and the Buratti--Horak--Rosa Conjecture}. 2024. \href{https://arxiv.org/abs/2402.08736}{arXiv:2402.08736}

\bibitem{ao24b} O. A\u{g}\i rseven, M. A. Ollis. \emph{A Coprime Buratti--Horak--Rosa Conjecture and Grid-Based Linear Realizations}. Revised January 2025. \href{https://arxiv.org/abs/2412.05750v2}{arXiv:2412.05750v2}

\bibitem{ao25} O. A\u{g}\i rseven, M. A. Ollis. \emph{Construction Techniques for Linear Realizations of Multisets with Small Support}. 2025. \href{https://arxiv.org/abs/2502.00164}{arXiv:2502.00164}

\end{thebibliography}
